\documentclass[11pt,a4paper]{article}
\usepackage{fontspec,amsmath,amssymb,amsthm,microtype}
\setmainfont{lmroman10-regular.otf}[BoldFont=lmroman10-bold.otf,ItalicFont=lmroman10-italic.otf,BoldItalicFont=lmroman10-bolditalic.otf]
\setmonofont{lmmono10-regular.otf}[Scale=MatchLowercase]
\usepackage[margin=25mm]{geometry}
\usepackage[unicode,hidelinks]{hyperref}
\setlength{\emergencystretch}{3em}
\setlength{\parskip}{0.4em}
\setlength{\parindent}{0pt}
\allowdisplaybreaks[2]
\raggedbottom
\title{\texorpdfstring{Push actual Wasserstein bounds through coordinate projections}{Push actual Wasserstein bounds through coordinate projections}}
\author{LangevinLean mathematical exposition}
\date{}
\begin{document}\maketitle
\noindent Source module: \texttt{Kinetic\allowbreak{}Langevin/\allowbreak{}Transport\allowbreak{}Map.\allowbreak{}lean}. The accompanying HTML pairs this mathematical exposition with the exact Lean statements and records the scope of the formalization.


These lemmas prove the Wasserstein effect of any actual measurable Lipschitz map between two normed spaces. They support coordinate projections and metric changes without assuming an optimal coupling.


\section{\texorpdfstring{Push each actual coupling through the map and then take the infimum}{Push each actual coupling through the map and then take the infimum}}


For any coupling c of the two input laws, apply the same measurable map T to both coordinates. The two marginal identities are proved by composing pushforwards with the original coupling marginals. The pointwise Lipschitz inequality and monotonicity of L\ensuremath{{}^{2}} norms bound the output coupling cost by q times the input cost, even though the input and output spaces may differ. Taking the infimum over every original coupling gives \(W_2(T_\#\mu,T_\#\nu)\le qW_2(\mu,\nu)\). A product coupling establishes that the input coupling type is nonempty. For the real-valued statement, the extended upper bound first implies finite output cost from finite initial cost, then the nonnegative-real coefficient and \texttt{to\allowbreak{}Real} identities give the claimed real inequality. This prevents an infinite transport cost from disappearing through conversion to a real number.



{\raggedright\noindent Lean declarations: \texttt{wasserstein2\_\allowbreak{}map\_\allowbreak{}lipschitz}, \texttt{wasserstein2\allowbreak{}Real\_\allowbreak{}map\_\allowbreak{}lipschitz}.\par}

\section{\texorpdfstring{Adjoin an independent common source before changing the output space}{Adjoin an independent common source before changing the output space}}


The same argument applies to a measurable map from an input state and a common independent innovation into a different normed output space. Extend every original coupling by the source law, prove both actual output marginals, and retain the unchanged input coupling cost. The pointwise contraction and L\ensuremath{{}^{2}} monotonicity bound each output cost. Taking the true infimum yields the cross-space random-map transport bound. In particular, adjoining one shared Gaussian velocity to two position laws has no cost of its own.



{\raggedright\noindent Lean declarations: \texttt{wasserstein2\_\allowbreak{}random\allowbreak{}Map\_\allowbreak{}lipschitz}.\par}

\end{document}
